% Einheit 2 von 4 – Winkel berechnen (bank/trigonometrie/e2.jsonl, zusammenbau v0.1)
%% TODO Zweigzeile Teil 1: was hier gelernt wird (Satz in Schülersprache aus dem Einheitstitel) – Einheit 2
\einheitenkopf[e2]{Einheit 2 von 4 · Winkel berechnen}
\zweigzeile{ab Kl. 10 (am Gymnasium ab Kl. 9) · P10 oft}
\setcounter{aufgabe}{18}

%% TODO Titel: Ich-kann-Satz fehlt in der Bank (Platzhalter: Kettenname) – Nr. 19
%% TODO Anweisung: ein Satz über der ersten Teilaufgabe fehlt in der Bank – Nr. 19
\begin{aufgabe}{Winkel berechnen}
\begin{teile}
\teil Gegeben: Hypotenuse $7$ cm und Gegenkathete $4$ cm. Gesucht: der Winkel $\alpha$. Was ist gesucht? \\ \kreuz{Seite: Gleichung umstellen} \\ \kreuz{Winkel: Umkehrtaste}
\end{teile}
%% TODO Darstellung für schwach fehlt in der Bank (trigonometrie-e2-k1-s1-v1; Abschnitt „Für schwache Schüler“)
\swz{Gegenkathete $3$ cm, Hypotenuse $8$ cm – wie groß ist $\alpha$?}{}
%% TODO Darstellung für schwach fehlt in der Bank (trigonometrie-e2-k1-s1-v2; Abschnitt „Für schwache Schüler“)
\swz{Gegenkathete $5$ cm, Hypotenuse $9$ cm – wie groß ist $\alpha$?}{}
%% TODO Darstellung für schwach fehlt in der Bank (trigonometrie-e2-k1-s1-v3; Abschnitt „Für schwache Schüler“)
\swz{Gegenkathete $2$ cm, Hypotenuse $5$ cm – wie groß ist $\alpha$?}{}
%% TODO Darstellung für schwach fehlt in der Bank (trigonometrie-e2-k1-s1-v4; Abschnitt „Für schwache Schüler“)
\swz{Gegenkathete $7$ cm, Hypotenuse $10$ cm – wie groß ist $\alpha$?}{}
%% TODO Darstellung für schwach fehlt in der Bank (trigonometrie-e2-k1-s1-v5; Abschnitt „Für schwache Schüler“)
\swz{Gegenkathete $5$ cm, Hypotenuse $8$ cm – wie groß ist $\alpha$?}{}
\swfrage{Was bleibt gleich, was ändert sich?}
%% TODO Darstellung für schwach fehlt in der Bank (trigonometrie-e2-k1-s2-v1; Abschnitt „Für schwache Schüler“)
\swz{Ankathete $5$ cm, Hypotenuse $7$ cm – wie groß ist $\alpha$?}{}
%% TODO Darstellung für schwach fehlt in der Bank (trigonometrie-e2-k1-s3-v1; Abschnitt „Für schwache Schüler“)
\swz{Gegenkathete $5$ cm, Ankathete $8$ cm – wie groß ist $\alpha$?}{}
\swz{Rechter Winkel bei $A$, $a = 10$ cm, $b = 3$ cm. Wie groß ist $\beta$?}{\dreieck{(0,0)}{(4,0)}{(0,3)}{a}{b}{c}{}{\beta}{}}
%% TODO Darstellung für schwach fehlt in der Bank (trigonometrie-e2-k1-s5-v1; Abschnitt „Für schwache Schüler“)
\swz{Gegenkathete von $\alpha$: $3$ cm, Hypotenuse $5$ cm. Berechne $\alpha$, dann $\beta$ als Ergänzung; kontrolliere $\beta$ mit dem Kosinus.}{}
%% TODO Darstellung für schwach fehlt in der Bank (trigonometrie-e2-k1-s6-v1; Abschnitt „Für schwache Schüler“)
\swz{Gegenkathete $2{,}3$ m, Hypotenuse $4{,}7$ m – wie groß ist $\alpha$, auf eine Dezimale?}{}
%% TODO Darstellung für schwach fehlt in der Bank (trigonometrie-e2-k1-s7-v1; Abschnitt „Für schwache Schüler“)
\swz{Eine Seilbahn überwindet $420$ m Höhe, das Seil ist $1\,300$ m lang. Wie groß ist der Steigungswinkel?}{}
\swz{Die Punkte $A$, $B$, $C$ bilden ein Dreieck mit rechtem Winkel bei $A$. Lies die Längen der Katheten ab und berechne den Winkel $\beta$ bei $B$.}{\begin{ksys}[xmin=-1,xmax=6,ymin=-1,ymax=5,ablesen] \punkt{1}{1}{A} \punkt{5}{1}{B} \punkt{1}{4}{C} \end{ksys}}
%% TODO Darstellung für schwach fehlt in der Bank (trigonometrie-e2-k1-s9-v1; Abschnitt „Für schwache Schüler“)
\swz{Im Dreieck $ABC$ ist die Höhe von $C$ auf $AB$ $4{,}5$ cm lang, $AC = 7{,}2$ cm. Wie groß ist $\alpha$ bei $A$?}{}
%% TODO Darstellung für schwach fehlt in der Bank (trigonometrie-e2-k1-s10-v1; Abschnitt „Für schwache Schüler“)
\swz[3]{Ankathete $5{,}2$ m, Hypotenuse $8$ m. Zeige rechnerisch, dass $\alpha$ rund $49^\circ$ groß ist.}{}
\swz[3]{Eine Seilbahn führt von der Talstation $A$ zur Bergstation $B$. Das Seil ist $610$ m lang, die waagerechte Strecke unter dem Seil $460$ m. Berechne den Winkel zwischen Seil und Waagerechter. \hfill (P10 2024 OS)}{}
\end{aufgabe}

%% TODO Titel: Ich-kann-Satz fehlt in der Bank (Platzhalter: Kettenname) – Nr. 20
%% TODO Anweisung: ein Satz über der ersten Teilaufgabe fehlt in der Bank – Nr. 20
\begin{aufgabe}{Winkel prüfen: die längere Kathete liegt dem größeren Winkel gegenüber}
%% TODO Darstellung für schwach fehlt in der Bank (trigonometrie-e2-k2-s1-v1; Abschnitt „Für schwache Schüler“)
\swz[3]{Die Katheten sind $4$ cm und $9$ cm lang. Tim sagt: Der Winkel gegenüber der $4$-cm-Kathete ist $66^\circ$. Kann das stimmen?}{}
\end{aufgabe}

%% TODO Titel: Ich-kann-Satz fehlt in der Bank (Platzhalter: Kettenname) – Nr. 21
%% TODO Anweisung: ein Satz über der ersten Teilaufgabe fehlt in der Bank – Nr. 21
\begin{aufgabe}{Umkehrtaste am Taschenrechner (SHIFT sin, cos, tan; Anzeige sin⁻¹)}
%% TODO Darstellung für schwach fehlt in der Bank (trigonometrie-e2-k3-s1-v1; Abschnitt „Für schwache Schüler“)
\swz{Mit der Umkehrtaste: $\mathrm{sin}^{-1}(0{,}42)$ – welcher Winkel, auf eine Dezimale?}{}
\end{aufgabe}

%% TODO Titel: Ich-kann-Satz fehlt in der Bank (Platzhalter: Kettenname) – Nr. 22
%% TODO Anweisung: ein Satz über der ersten Teilaufgabe fehlt in der Bank – Nr. 22
\begin{aufgabe}{Winkel berechnen – Fehler finden, Begründen, Anwendung}
\swz[3]{Gegenkathete von $\alpha$: $3$ cm, Ankathete: $8$ cm. Wo ist der Fehler? \rechnung{\mathrm{tan}\,\alpha &= \frac{8}{3} \\ \alpha &\approx 69{,}4^\circ}}{}
\swfrage{Ole tippt $\mathrm{sin}^{-1}(1{,}3)$ und der Taschenrechner meldet einen Fehler. Warum?}
\swz[3]{Eine Rollstuhlrampe darf höchstens $3{,}4^\circ$ steil sein. Eine $5$ m lange Rampe überwindet $0{,}32$ m. Ist sie zulässig?}{}
\end{aufgabe}

\uebersichtskasten{Winkel aus zwei Seiten: Bilde das Verhältnis der beiden gegebenen Seiten vom gesuchten Winkel aus (Gegenkathete oben, Hypotenuse oder Ankathete unten) und drücke die Umkehrtaste – sin⁻¹, cos⁻¹ oder tan⁻¹ macht aus dem Verhältnis den Winkel. \\ Hypotenuse 11 cm, Gegenkathete von α 7 cm: sin α = 7/11 ≈ 0,636 \quad α = sin⁻¹(7/11) ≈ 39,5° \\ Derselbe Bruch mit der Ankathete: cos β = 7/11 \quad β ≈ 50,5°  (α und β ergänzen sich zum rechten Winkel) \\ Beide Katheten 7 cm und 9 cm: tan α = 7/9 \quad α = tan⁻¹(7/9) ≈ 37,9° \\ Merkregel: Seite gesucht → Gleichung umstellen. Winkel gesucht → Verhältnis, dann Umkehrtaste; der Sinus- und der Kosinuswert sind nie größer als eins.}
